%!TEX root = rapport.tex

\chapter{Introduction}

We will in this project develop an automatic method for visualization of 
$\beta$-reduction graphs of lambda terms. The method has been implemented in 
the programming language Python.

Until now there did not exist any software for reduction graph visualization.
Researchers wanting to visualize reduction graphs have therefore been limited
to either their own drawings or to what they can imagine. Since reduction
graphs can be very large, this posed a very tight constraint on what could be
visually presented. 

In our opinion, software for this visualization task presents two primary
challenges: The drawn graphs must be \emph{informative} and
\emph{good-looking}. These two matters are not entirely independent of each
other; very often the particular layout of a graphs two-dimensional embedding
greatly reduces clutter and ``noise'', thus adding information. The opposite
can of course also be the case, which is why the two problems are
interchanged.

The intended user of the software is a theoretical computer scientist that
does research in the field of lambda calculus. Our hope is that researchers
can use this tool to gain new insight into the lambda calculus, and that it
will serve as a source of inspiration for new research because it offers a new
look at a well-known concept.

% \section{Project constraints \& risks}
% Because reduction graphs very often tend to become huge, there arises the
% problem of deci ding whether they can be drawn on a plane at all. We will
% perhaps have to develop a routine for deciding this, along with routines for
% reordering the graph. Alternatively, a set of ``best practices'' for drawing
% the nodes could be developed that would minimize the amount of crossing edges
% as the film progresses forward. Yet more alternatively, methods for 3D-drawing
% could be considered. However, this will demand great care when determining the
% look of the images, as a 3-dimensional rendering of a large graph very easily
% could inhibit understanding and confuse the beholder, which is the opposite of
% what we are working towards.
% 
% Furthermore, there is a design choice to be made whether each frame of the
% film should be drawn immediately after the reduction step has been completed,
% making the film rendering and the term reductions synchronized, or whether the
% rendering should begin only after all desired (or possible) reductions have
% been computed.
